\documentclass[12pt,a4paper]{article}

% ─── Codificación y lenguaje ─────────────────────────────────────────────────
\usepackage[utf8]{inputenc}
\usepackage[T1]{fontenc}
\usepackage[spanish,es-noquoting]{babel}

% ─── Diseño de página ────────────────────────────────────────────────────────
\usepackage[top=2.4cm, bottom=2.8cm, left=2.8cm, right=2.8cm]{geometry}
\usepackage{parskip}
\usepackage{microtype}
\usepackage{fancyhdr}
\pagestyle{fancy}
\fancyhf{}
\fancyhead[L]{\small\textit{Actividad formativa --- Minutas IA}}
\fancyhead[R]{\small\textit{Ingeniería Informática}}
\fancyfoot[C]{\thepage}
\renewcommand{\headrulewidth}{0.4pt}

% ─── Tipografía ──────────────────────────────────────────────────────────────
\usepackage{lmodern}
\usepackage{setspace}
\setstretch{1.12}

% ─── Colores y cajas ─────────────────────────────────────────────────────────
\usepackage{xcolor}
\definecolor{azulprimario}{RGB}{20, 70, 140}
\definecolor{azulclaro}{RGB}{230, 240, 255}
\definecolor{verdeok}{RGB}{10, 110, 50}
\definecolor{verdeclaro}{RGB}{230, 248, 236}
\definecolor{rojowarn}{RGB}{160, 30, 20}
\definecolor{rojoclaro}{RGB}{255, 235, 230}
\definecolor{naranjainf}{RGB}{160, 90, 0}
\definecolor{naranjaclaro}{RGB}{255, 245, 220}
\definecolor{backcolour}{rgb}{0.96,0.96,0.94}
\definecolor{codegreen}{rgb}{0,0.52,0}
\definecolor{codegray}{rgb}{0.5,0.5,0.5}
\definecolor{codepurple}{rgb}{0.5,0,0.62}
\definecolor{codeblue}{rgb}{0.0,0.18,0.72}

\usepackage{tcolorbox}
\tcbuselibrary{skins, breakable}

\newtcolorbox{conceptbox}[1]{
  title={\textbf{#1}},
  colback=azulclaro, colframe=azulprimario,
  fonttitle=\color{white}, coltitle=white,
  colbacktitle=azulprimario,
  attach boxed title to top left={yshift=-2mm, xshift=4mm},
  boxed title style={rounded corners, size=small},
  rounded corners, breakable, left=6pt, right=6pt,
  top=8pt, bottom=6pt,
}

\newtcolorbox{warnbox}[1]{
  title={\textbf{#1}},
  colback=rojoclaro, colframe=rojowarn,
  fonttitle=\color{white}, coltitle=white,
  colbacktitle=rojowarn,
  attach boxed title to top left={yshift=-2mm, xshift=4mm},
  boxed title style={rounded corners, size=small},
  rounded corners, breakable, left=6pt, right=6pt,
  top=8pt, bottom=6pt,
}

\newtcolorbox{examplebox}[1]{
  title={\textbf{#1}},
  colback=verdeclaro, colframe=verdeok,
  fonttitle=\color{white}, coltitle=white,
  colbacktitle=verdeok,
  attach boxed title to top left={yshift=-2mm, xshift=4mm},
  boxed title style={rounded corners, size=small},
  rounded corners, breakable, left=6pt, right=6pt,
  top=8pt, bottom=6pt,
}

\newtcolorbox{infobox}[1]{
  title={\textbf{#1}},
  colback=naranjaclaro, colframe=naranjainf,
  fonttitle=\color{white}, coltitle=white,
  colbacktitle=naranjainf,
  attach boxed title to top left={yshift=-2mm, xshift=4mm},
  boxed title style={rounded corners, size=small},
  rounded corners, breakable, left=6pt, right=6pt,
  top=8pt, bottom=6pt,
}

\newtcolorbox{experbox}[1]{
  title={\textbf{#1}},
  colback=white, colframe=azulprimario!60,
  fonttitle=\color{azulprimario},
  attach boxed title to top left={yshift=-2mm, xshift=4mm},
  boxed title style={rounded corners, size=small},
  rounded corners, breakable, left=6pt, right=6pt,
  top=8pt, bottom=6pt,
}

% ─── Código y tablas ─────────────────────────────────────────────────────────
\usepackage{listings}
\usepackage{booktabs}
\usepackage{amsmath}
\usepackage{enumitem}
\usepackage{array}
\usepackage[hidelinks, colorlinks=true,
            linkcolor=azulprimario,
            urlcolor=azulprimario]{hyperref}

\lstdefinestyle{shellstyle}{
  backgroundcolor=\color{backcolour},
  basicstyle=\ttfamily\small,
  breaklines=true, frame=single,
  rulecolor=\color{black!25},
  showstringspaces=false, xleftmargin=14pt,
}

\lstdefinestyle{envstyle}{
  backgroundcolor=\color{backcolour},
  basicstyle=\ttfamily\small,
  breaklines=true, frame=single,
  rulecolor=\color{black!25},
  showstringspaces=false, xleftmargin=14pt,
}

\usepackage{titlesec}
\titleformat{\section}
  {\Large\bfseries\color{azulprimario}}{\thesection.}{0.6em}{}[\vspace{-0.4em}\rule{\linewidth}{0.5pt}]
\titleformat{\subsection}
  {\large\bfseries\color{azulprimario!80!black}}{\thesubsection.}{0.5em}{}

\newcommand{\checkbox}{\fbox{\rule{0pt}{1.5ex}\rule{1.5ex}{0pt}}}

% ─── Documento ─────────────────────────────────────────────────────────────────
\title{
  \vspace{-0.8cm}
  {\large Ingeniería Informática --- Actividad formativa}\\[0.4em]
  {\LARGE\bfseries\color{azulprimario}
    Experimentación guiada con Minutas IA}\\[0.3em]
  {\normalsize LLM \textbullet{} RAG \textbullet{} Redis \textbullet{}
   LangChain \textbullet{} Ollama}
}
\author{IA Embebida en Sistemas Computacionales}
\date{Semestre I --- 2026}

\begin{document}

\maketitle
\thispagestyle{fancy}

\begin{infobox}{Ficha de la actividad}
\begin{tabular}{@{}ll@{}}
\toprule
\textbf{Asignatura} & IA Embebida en Sistemas Computacionales \\
\textbf{Duración} & 50--60 minutos \\
\textbf{Modalidad} & Individual o parejas \\
\textbf{Tipo} & Formativa (entregable breve, sin nota numérica) \\
\bottomrule
\end{tabular}
\end{infobox}

\section{Propósito}

Profundizar en la arquitectura entregada (LLM + RAG + Redis + LangChain) mediante
\textbf{experimentación controlada}, considerando que \textbf{no todos los equipos
tienen la misma capacidad} para ejecutar modelos de IA locales.

Al finalizar deberías poder explicar, con evidencia de tu propia prueba:

\begin{enumerate}[noitemsep]
  \item Qué aporta el \textbf{RAG} frente a un LLM aislado.
  \item Qué aporta \textbf{Redis} como memoria de sesión.
  \item Cómo el \textbf{pre-prompt} y el \textbf{tamaño del modelo} afectan
        calidad, latencia y consumo de recursos.
  \item Qué \textit{trade-offs} harías al desplegar este sistema en un jardín
        infantil real.
\end{enumerate}

\section{Antes de empezar}

\subsection{Requisitos mínimos}

\begin{itemize}[noitemsep]
  \item Docker Desktop o Docker Engine + Compose v2.
  \item Conexión a internet (solo la primera vez, para imágenes y modelos).
  \item Repositorio clonado con \texttt{make up} funcional, o acceso al servidor
        del curso (Perfil R).
\end{itemize}

\subsection{Panel de trabajo}

Abre \url{http://localhost:3000} y ubica:

\begin{itemize}[noitemsep]
  \item Formulario de generación de minutas.
  \item Bloque \textbf{Pre-prompt y vectorización RAG} (parte inferior).
  \item Presets de actividades (María, Pedro consulta 1 y 2).
\end{itemize}

\section{Paso 0: Elige tu perfil de hardware (5 min)}

No todos corren el mismo modelo. Elige un perfil según tu equipo \textbf{antes}
de los experimentos.

\begin{table}[h]
\centering
\caption{Perfiles de hardware para el curso}
\begin{tabular}{@{}llll@{}}
\toprule
\textbf{Perfil} & \textbf{RAM aprox.} & \textbf{Modelo LLM} & \textbf{Comando} \\
\midrule
L --- Ligero & 8 GB, sin GPU & \texttt{llama3.2:1b} & \texttt{make perfil-ligero} \\
M --- Medio  & 16 GB & \texttt{llama3.2} (3B) & \texttt{make perfil-medio} \\
A --- Alto   & 16+ GB + GPU / Apple Silicon & \texttt{llama3.1:8b} & \texttt{make perfil-alto} \\
R --- Remoto & PC débil o sin Docker & Ollama del laboratorio & Sección~\ref{sec:perfil-r} \\
\bottomrule
\end{tabular}
\end{table}

Después de elegir perfil:

\begin{lstlisting}[style=shellstyle]
make perfil-ligero   # o perfil-medio / perfil-alto
docker compose down
make up
make ready
\end{lstlisting}

Anota en tu bitácora:

\begin{itemize}[noitemsep]
  \item Perfil elegido: \rule{6cm}{0.4pt}
  \item Modelo activo (panel «Estado del sistema» o \texttt{make ready}):
        \rule{4cm}{0.4pt}
  \item Tiempo de la \textbf{primera} respuesta del LLM (s):
        \rule{2cm}{0.4pt}
\end{itemize}

\begin{warnbox}{Nota docente}
El objetivo no es comparar notas por velocidad, sino \textbf{razonar sobre
trade-offs} con el hardware disponible.
\end{warnbox}

\section{Experimento 1: ¿Qué aporta el RAG? (12 min)}

\begin{experbox}{Experimento 1 --- RAG}
\subsection*{Hipótesis}
\textit{Sin contexto documental, el LLM ignora la normativa JUNJI/MINSAL y responde
de forma genérica.}

\subsection*{Procedimiento}
\begin{enumerate}[noitemsep]
  \item En el panel inferior, revisa el \textbf{pre-prompt} y al menos
        \textbf{2 fragmentos} del corpus indexado.
  \item Usa el preset \textbf{«Actividad 1 --- María (lactosa)»}.
  \item Pulsa \textbf{«Vista previa RAG»} y observa el embedding de la consulta
        y los fragmentos recuperados (\texttt{similarity\_score}).
  \item Genera la minuta y compara el resultado con los fragmentos recuperados.
  \item Abre \texttt{documentos\_normativos/normativa\_junji\_demo.txt} y verifica
        si algún fragmento menciona: kcal 200--350, alérgenos o ingredientes chilenos.
\end{enumerate}

\subsection*{Registro (completa en tu entrega)}
\begin{table}[h]
\centering
\small
\begin{tabular}{@{}p{7.2cm}p{6.5cm}@{}}
\toprule
\textbf{Pregunta} & \textbf{Tu observación} \\
\midrule
¿Cuántos fragmentos recuperó el retriever? & \\[1.8ex]
¿Algún fragmento menciona kcal o alérgenos? Cita una línea. & \\[1.8ex]
¿El plato generado respeta la lactosa? & \\[1.8ex]
¿El campo \texttt{fuente\_normativa} aparece? ¿Es creíble? & \\
\bottomrule
\end{tabular}
\end{table}

\subsection*{Pregunta de reflexión}
Si mañana se publica una nueva circular JUNJI, \textbf{¿qué componente actualizas
sin reentrenar el LLM?} Justifica en 2--3 líneas.
\end{experbox}

\section{Experimento 2: Memoria con Redis (10 min)}

\begin{experbox}{Experimento 2 --- Redis}
\subsection*{Hipótesis}
\textit{Con el mismo \texttt{session\_id}, el sistema no debería obligar a repetir
las alergias en cada consulta.}

\subsection*{Procedimiento}
\begin{enumerate}[noitemsep]
  \item Preset \textbf{«Actividad 2 --- Pedro consulta 1»} $\rightarrow$ Generar minuta.
  \item Preset \textbf{«Actividad 2 --- Pedro consulta 2 (memoria)»} $\rightarrow$
        Generar minuta \textbf{sin} editar el \texttt{session\_id}.
  \item (Opcional) Repite la consulta 2 con \texttt{session\_id} =
        \texttt{test:pedro-nuevo} y compara.
\end{enumerate}

\subsection*{Verificación técnica (elige una)}
\begin{lstlisting}[style=shellstyle, caption={Opción A: Redis CLI}]
# Opcion A: Redis CLI
docker compose exec redis redis-cli LRANGE minuta:test:pedro 0 -1

# Opcion B: observacion en el panel de historial (UI)
\end{lstlisting}

\subsection*{Registro}
\begin{table}[h]
\centering
\small
\begin{tabular}{@{}p{5.5cm}cc@{}}
\toprule
\textbf{Escenario} & \textbf{¿Excluyó huevo/mariscos?} & \textbf{¿Por qué?} \\
\midrule
Consulta 2, mismo \texttt{session\_id} & & \\[1.5ex]
Consulta 2, \texttt{session\_id} distinto (opc.) & & \\
\bottomrule
\end{tabular}
\end{table}

\subsection*{Pregunta de reflexión}
¿Redis almacena «inteligencia» o solo \textbf{historial de mensajes}? ¿Qué pasaría
si el TTL pasara de 24 h a 7 días en un jardín con 80 niños?
\end{experbox}

\section{Experimento 3: Pre-prompt y calidad del JSON (10 min)}

\begin{experbox}{Experimento 3 --- Pre-prompt y JSON}
\subsection*{Procedimiento}
\begin{enumerate}[noitemsep]
  \item Lee el pre-prompt en el panel inferior (reglas 1--9).
  \item Genera una minuta con esta consulta:
        \begin{quote}
          \textit{«Niño de 18 meses, sin alergias. Once del miércoles con
          ingredientes de temporada.»}
        \end{quote}
  \item Abre \textbf{«Ver JSON completo»} en el resultado.
  \item Marca con \checkbox{} (sí) o vacío (no):
\end{enumerate}

\begin{itemize}[noitemsep]
  \item[\checkbox] \texttt{nino\_id} es un valor real (no \texttt{"string"})
  \item[\checkbox] \texttt{generada\_en} es fecha ISO válida
  \item[\checkbox] \texttt{kcal} está entre 200 y 350
  \item[\checkbox] \texttt{verificado\_alergenos} es coherente con la consulta
  \item[\checkbox] Ingredientes con nombres locales (porotos, choclo, etc.)
\end{itemize}

\subsection*{Variante según perfil de hardware}
\begin{table}[h]
\centering
\small
\begin{tabular}{@{}lp{9.5cm}@{}}
\toprule
\textbf{Perfil} & \textbf{Variante extra} \\
\midrule
L & Si tarda $>$90 s o falla: \texttt{RETRIEVER\_K=1} en \texttt{.env}, reinicia API. \\
M & Sin variante; registra tiempo de respuesta. \\
A & Prueba \texttt{LLM\_TEMPERATURE=0.5} y compara creatividad vs.\ consistencia. \\
R & Mismo flujo en servidor remoto; anota latencia de red. \\
\bottomrule
\end{tabular}
\end{table}

\subsection*{Pregunta de reflexión}
¿Basta con \texttt{format=json} en Ollama para garantizar un JSON \textbf{válido
para producción}? Menciona al menos un caso que Pydantic podría rechazar.
\end{experbox}

\section{Experimento 4: Trade-off de modelos (8 min)}

\begin{experbox}{Experimento 4 --- Trade-offs}
\textbf{Objetivo:} relacionar potencia del equipo con calidad y latencia, sin
competir por hardware.

\subsection*{Procedimiento}
\begin{enumerate}[noitemsep]
  \item Anota tu perfil y modelo actual.
  \item Si tu equipo y el docente lo autorizan, cambia de perfil (ej.\ M $\rightarrow$ L)
        y repite la consulta de María.
  \item Completa la tabla (si no puedes cambiar de perfil, completa solo tu fila):
\end{enumerate}

\begin{table}[h]
\centering
\small
\begin{tabular}{@{}lcc@{}}
\toprule
\textbf{Métrica} & \textbf{Mi perfil} & \textbf{Otro perfil (opc.)} \\
\midrule
Modelo & & \\[1.2ex]
Tiempo 1.\textsuperscript{a} respuesta (s) & & \\[1.2ex]
¿JSON válido a la 1.\textsuperscript{a}? & & \\[1.2ex]
¿Plato nutricionalmente plausible? & & \\[1.2ex]
¿Alucinó ingredientes raros? & & \\
\bottomrule
\end{tabular}
\end{table}

\subsection*{Pregunta de reflexión}
En producción, ¿siempre conviene el modelo más grande? Nombra \textbf{dos criterios}
además de la calidad del texto (costo, latencia, privacidad, energía, etc.).
\end{experbox}

\section{Entregable formativo (máx.\ 1 página)}

Sube un PDF con:

\begin{enumerate}[noitemsep]
  \item \textbf{Perfil de hardware} usado y modelo Ollama activo.
  \item \textbf{Tablas de registro} de los experimentos 1 y 2 (aunque haya fallos).
  \item \textbf{Respuestas breves} (3--5 líneas) a \textbf{dos} preguntas de reflexión
        a elección.
  \item \textbf{Captura} del panel «Pre-prompt y vectorización RAG» con al menos un
        fragmento recuperado.
  \item \textbf{Un problema real} detectado y \textbf{qué capa} falló: LLM, RAG,
        Redis, prompt o validación.
\end{enumerate}

\begin{examplebox}{Si no pudiste ejecutar localmente}
Indica uso del \textbf{Perfil R} y adjunta capturas del servidor del curso con tu
\texttt{session\_id} identificable.
\end{examplebox}

\section{Perfil R --- Estudiantes sin PC potente}
\label{sec:perfil-r}

Si tu laptop no puede ejecutar Ollama de forma razonable:

\begin{enumerate}[noitemsep]
  \item Solicita al docente la URL del servidor (ej.\
        \texttt{http://lab-ia.universidad.cl:11434}).
  \item Crea un archivo \texttt{.env} local:
\end{enumerate}

\begin{lstlisting}[style=envstyle]
OLLAMA_BASE_URL=http://<servidor-del-curso>:11434
LLM_MODEL=llama3.2:1b
LLM_NUM_CTX=2048
RETRIEVER_K=2
\end{lstlisting}

\begin{enumerate}[noitemsep]
  \setcounter{enumi}{2}
  \item Levanta API + Redis + Web según indique el docente, o usa el frontend del
        laboratorio.
  \item En la entrega documenta: \textbf{«Ejecución remota --- sin inferencia local»}.
\end{enumerate}

\subsection{Alternativa mínima sin Docker}

Observa la demo del docente y completa los experimentos 1 y 3 \textbf{teóricamente},
citando capturas del proyector y \texttt{normativa\_junji\_demo.txt}. El docente
puede aceptar esta vía solo si quedó sin equipo; prioriza siempre el Perfil L.

\section{Criterios de logro (autoevaluación)}

\begin{table}[h]
\centering
\begin{tabular}{@{}lc@{}}
\toprule
\textbf{Criterio} & \textbf{Lo logré} \\
\midrule
Identifiqué qué hace el RAG en la arquitectura & \checkbox \\[0.8ex]
Demostré memoria de sesión con Redis & \checkbox \\[0.8ex]
Interpreté el pre-prompt y los fragmentos vectorizados & \checkbox \\[0.8ex]
Relacioné tamaño de modelo con latencia/calidad & \checkbox \\[0.8ex]
Distinguí fallo de LLM vs.\ validación/RAG & \checkbox \\
\bottomrule
\end{tabular}
\end{table}

\section{Para seguir explorando (opcional)}

\begin{itemize}[noitemsep]
  \item Agrega un párrafo a \texttt{documentos\_normativos/} y ejecuta
        \textbf{Reindexar RAG}.
  \item Modifica una regla en \texttt{app/system\_prompt.py}, reconstruye la API
        y observa el cambio.
  \item Compara \texttt{RETRIEVER\_K=2} vs.\ \texttt{RETRIEVER\_K=6} y mide latencia.
\end{itemize}

\section{Comandos de referencia rápida}

\begin{lstlisting}[style=shellstyle]
make perfil-ligero    # 8 GB RAM
make perfil-medio     # 16 GB RAM (default clase)
make perfil-alto      # equipos potentes
make up && make ready
make web-dev          # frontend en desarrollo
docker compose logs -f api
\end{lstlisting}

\vfill
\begin{center}
\small
\textit{Actividad alineada con el caso de estudio: sistema de minutas alimentarias
JUNJI --- Ingeniería Informática.}
\end{center}

\end{document}
